\documentclass[10pt,a4paper]{amsart}
\usepackage[margin=19mm]{geometry}
\usepackage{amsmath,amssymb,mathtools}
\usepackage[scr=rsfs,cal=euler]{mathalfa}
\usepackage{xfrac}
\usepackage[unicode,hidelinks,bookmarks=false]{hyperref}
\newcommand{\ie}{i.e.\ }
\newcommand{\cf}{cf.\ }
\newcommand{\resp}{respectively\ }
\newcommand{\Subsection}{\subsection}
\providecommand{\nobreakdash}{\nobreak\hbox{-}}
\setlength{\emergencystretch}{2em}
\multlinegap0pt
\usepackage{amssymb}
\usepackage{mathtools}
\newtheorem{definition}{Definition}
\newtheorem{lemma}{Lemma}
\newtheorem{corollary}{Corollary}
\newcommand{\D}{\mathcal{D}}
\DeclareMathOperator{\ddell}{ddell}
\DeclareMathOperator{\dell}{dell}
\DeclareMathOperator{\modd}{mod}
\newcommand{\oemb}{\stackrel{\oplus}{\hookrightarrow}}
\newcommand{\syz}{\Omega_{\mathcal{D}}}
\title{Degree-shift derived invariance of derived delooping levels}
\author{Jiaqun Wei, Kaili Wu, Weiqing Cao}
\date{Source: arXiv:2608.29634v1 (CC BY 4.0)}
\begin{document}
\maketitle

\noindent\emph{Source and licence.} Degree-shift derived invariance of derived delooping levels. Version arXiv:2608.29634v1 (CC BY 4.0), distributed by arXiv under the Creative Commons Attribution 4.0 International licence, http://creativecommons.org/licenses/by/4.0/, as recorded in the arXiv licence metadata for this exact version and shown on the abstract page https://arxiv.org/abs/2608.29634v1. Full licence notice: \texttt{licenses/arxiv-2608.29634-CC-BY-4.0.txt}.\par\smallskip

\noindent\textit{Editorial scope.} Counted unit: the article's lemma lem:transfer (source0.tex lines 602--612). The article's complete original proof of it is delivered with it (source0.tex lines 614--671, one \texttt{proof} environment of 58 lines). No mathematics is written, repaired or re-derived: the statement and the proof are the article's own lines, in the article's own notation, and no retained line is substituted or re-flowed.  Retained uncounted as the source-local context the proof needs: lines 325--331; lines 379--391; lines 440--450; lines 520--530; lines 582--592.  Nothing was omitted from any retained span, so the package carries no omission record.  No retained line contains a citation, so the wrapper carries no bibliography and no bibliography entry.  No retained line contains a figure environment, an image-inclusion command or drawing code, so no image is delivered and the package declares no asset dependency.  Editorial formatting only, in addition: an AMS \texttt{amsart} preamble that restates the article's own declarations and macros used by the retained lines, namely the environments \texttt{definition}, \texttt{lemma}, \texttt{corollary} on separate counters, together with the standard packages those lines need.  The retained environments print in this document as Definition 1, Lemma 1, Corollary 1 in order of appearance, whereas the article prints the same items with the numbering of its own full document.  The complete native source of the article as distributed in the arXiv e-print for this exact version is bundled unchanged as \texttt{sources/208871/source0.tex}.
\begin{definition}\label{def:k-dell}
Let $M \in \modd R$. The \emph{$k$-delooping level} of $M$ is
\[
k\text{-}\dell(M) = \inf\{n \geq 0 \mid \Omega^n(M)  \oemb  \Omega^{n+k}(N) \text{ for some } N \in \modd R\}.
\]
If no such $n$ exists, we set $k\text{-}\dell(M) = \infty$. The \emph{$k$-delooping level} of $R$ is $k\text{-}\dell(R) = \sup\{k\text{-}\dell(S) \mid S \text{ simple}\}$.
\end{definition}

\begin{definition}\label{def:classical-ddell}
$(1)$ We call the integer $k$ in Definition \ref{def:k-dell} the \emph{delooping degree}. Thus, $k\text{-}\dell(M)$ means the delooping level of $M$ with degree $k$.
 
$(2)$ Let $M \in \modd R$. We say that $M$ admits a \emph{derived delooping resolution $($of length $n$$)$ with degree $k$}, or shortly \emph{$k$-$\ddell$ resolution}, if there is a finite exact sequence in $\modd R$:
\[
0 \longrightarrow D_n \longrightarrow D_{n-1} \longrightarrow \cdots \longrightarrow D_1 \longrightarrow D_0 \longrightarrow M \longrightarrow 0
\]
such that for each $i \in \{0, 1, \ldots, n\}$, there exists a module $N_i \in \modd R$ with
\[
\Omega^{n-i}(D_i) \oemb \Omega^{n+k}(N_i),
\]
meaning that $\Omega^{n-i}(D_i)$ is a direct summand of $\Omega^{n+k}(N_i)$, i.e., $(i+k)$-$\dell(D_i)\le n-i$.
\end{definition}

\begin{lemma}\label{lem:image-bounded}
Let $a \leq b$. Then 

$(1)$ $ \mathcal{F}(\D_{[a,b]}(R)) \subseteq \D_{[a,b+k_T]}(S)$.
In particular, $\mathcal{F}(\modd R) \subseteq \D_{[0,k_T]}(S)$.

$(2)$ \(
\mathcal G(\D_{[a,b]}(S)) \subseteq \D_{[a-k_T,\, b]}(R).
\)
In particular, \(\mathcal G(\mathrm{mod}\,S)\subseteq \D_{[-k_T,0]}(R)\).
\end{lemma}

\begin{corollary} \label{cor:iterative-syzygy}
$(1)$
\(
\syz^j(\mathcal{G}(\syz^i(C))) \simeq \syz^{j+i}(\mathcal{G}(C)) 
\) holds for $C \in \modd S$ and $i,j \geq 0$.

$(2)$ 
\(
\syz^j(\mathcal{F}(\syz^i(C))) \simeq \syz^{j+i}(\mathcal{F}(C))
\) holds for $C \in \modd R$, $j \geq k_T$ and $i\ge 0$.
\end{corollary}

\begin{corollary}\label{cor:key-syzygy}
$(1)$ If $X \in \modd R$, then
\(
\syz^{0}(\mathcal{G}(\syz^{k_T}(\mathcal{F}(X)))) \simeq \syz^{k_T}(X).
\)

$(2)$ If $X \in \modd S$, then
\(
\syz^{k_T}(\mathcal{F}(\syz^{0}(\mathcal{G}(X)))) \simeq \syz^{k_T}(X).
\)
\end{corollary}

\begin{lemma} \label{lem:transfer}
$(1)$ Let \(N\in\modd R\) and set $
M_N:=\syz^{k_T}(\mathcal{F}(N))\in\modd S.
$
If \(M_N\) admits a \(k\)-\(\ddell\) resolution of length \(n\), then \(\syz^{k_T}(N)\) admits a \(k\)-\(\ddell\) resolution of length \(n\).

$(2)$ Let \(M\in\modd S\) and set $
N_M:=\syz^{0}(\mathcal{G}(M))\in\modd R.
$
If \(N_M\) admits a \(k\)-\(\ddell\) resolution of length \(n\), then \(\syz^{k_T}(M)\) admits a \(k\)-\(\ddell\) resolution of length \(n\).
\end{lemma}

\begin{proof} (1)
It is clear that $M_N\in\modd S$ since $\mathcal{F}(N)\in\D_{[0,k_T]}(S)$ by Lemma \ref{lem:image-bounded}.
Suppose $M_N$ has a classical $k$-$\ddell$ resolution of length \(n\) in $\modd S$:
\[
0 \longrightarrow D_n \longrightarrow D_{n-1} \longrightarrow \cdots \longrightarrow D_0 \longrightarrow M_N \longrightarrow 0. \tag{$\dagger$}
\]
By Definition \ref{def:classical-ddell}, for each \(i=0,\ldots,n\), there exists \(L_i\in\modd S\) such that
\[
\Omega^{n-i}(D_i) \oemb \Omega^{n+k}(L_i). \tag{$\ddagger$}
\]
That is, there are $Y_i$'s such that $\Omega^{n-i}(D_i)\oplus Y_i\simeq \Omega^{n+k}(L_i)$ in $\modd S$. The exact sequence consists of short exact sequences in $\modd S$ 
\[
    0\to C_{i+1}\to D_i\to C_i\to 0, 
\]
where $C_n=D_n$ and $C_0=M_N$.

Applying $\mathcal{G}$ to these exact sequences gives triangles in $\mathcal{G}(\modd S)\subseteq \D_{[-k_T,0]}(R)$:
 \[
    \mathcal{G}(C_{i+1})\to \mathcal{G}(D_i)\to \mathcal{G}(C_i)\to . 
\]
 
Now applying $\syz^{0}$, we obtain triangles in $\modd R$, and hence exact sequences in $\modd R$:
 \[
    0\to \syz^{0}(\mathcal{G}(C_{i+1}))\to \syz^{0}(\mathcal{G}(D_i))\to \syz^{0}(\mathcal{G}(C_i))\to 0,
\]
which combine to give an exact sequence in $\modd R$:
\[
0 \longrightarrow \syz^{0}(\mathcal{G}(D_n)) \longrightarrow \cdots \longrightarrow \syz^{0}(\mathcal{G}(D_0)) \longrightarrow \syz^{0}(\mathcal{G}(M_N)) \longrightarrow 0. \tag{$\star$}
\]

 Moreover, from \((\ddagger)\), applying $\mathcal{G}$ gives an embedding
\[
\mathcal{G}(\Omega^{n-i}(D_i)) \oemb \mathcal{G}(\Omega^{n+k}(L_i)).
\]
By Corollary \ref{cor:iterative-syzygy} (1), we obtain
\[
\syz^{0}(\mathcal{G}(\Omega^{n-i}(D_i))) \simeq \Omega^{n-i}\bigl(\mathcal{G}(D_i)\bigr) \simeq \Omega^{n-i}\bigl(\syz^{0}(\mathcal{G}(D_i))\bigr),
\]
and similarly,
\[
\syz^{0}(\mathcal{G}(\Omega^{n+k}(L_i))) \simeq \Omega^{n+k}\bigl(\mathcal{G}(L_i)\bigr) \simeq \Omega^{n+k}\bigl(\syz^{0}(\mathcal{G}(L_i))\bigr).
\]
Combining these yields
\[
\Omega^{n-i}\bigl(\syz^{0}(\mathcal{G}(D_i))\bigr) \oemb \Omega^{n+k}\bigl(\syz^{0}(\mathcal{G}(L_i))\bigr),
\]
which is precisely the required condition for \((\star)\) to be a $k$-$\ddell$ resolution in $\modd R$ (note that all $\syz^{0}(\mathcal{G}(L_i))\in \modd R$).

Since
\[
\syz^{0}(\mathcal{G}(M_N))
= \syz^{0}(\mathcal{G}(\syz^{k_T}(\mathcal{F}(N))))
\simeq \syz^{k_T}(N)
\] 
by Corollary \ref{cor:key-syzygy} (1), we see that \((\star)\) is a $k$-$\ddell$ resolution of length $n$ of $\syz^{k_T}(N)$.

(2) The proof is dual to (1).
\end{proof}

\end{document}
